\documentclass[12pt,a4paper,oneside]{article}
\usepackage[brazil]{babel}
\usepackage[utf8]{inputenc}
\usepackage{olymp}


\newlength{\exmpwidouf}
\exmpwidinf=10cm
\exmpwidouf=6cm


\setcounter{problem}{1}

\begin{document}

\begin{problem}{Número de Euler}{euler}{2}

O número de Euler, $e$, que recebeu este nome em homenagem ao matemático suíço Leonhard Euler, é a base dos logaritmos naturais. Ele pode ser definido por:

$$e=\sum_{n=0}^\infty{\frac{1}{n!}}=\frac{1}{0!}+\frac{1}{1!}+\frac{1}{2!}+\frac{1}{3!}+\frac{1}{4!}+\dots$$

onde $n!$ é o fatorial de $n$, dado por:

$$n! = 1\times 2 \times \dots \times (n-1) \times n$$

obs.: $0!=1$.

Faça um programa para calcular uma aproximação do número de Euler com $M$ termos, isto é:

$$e \approx \sum_{n=0}^{M-1}{\frac{1}{n!}}=\frac{1}{0!}+\frac{1}{1!}+\frac{1}{2!}+\dots+\frac{1}{(M-1)!}$$

\Notes

Você pode criar uma função para o cálculo do fatorial.

\InputFile

A entrada é composta por um único valor inteiro $M$, indicando o número de termos a serem utilizados para a aproximação. Restrição:
$1 \leq M \leq 10$.

\OutputFile

Seu programa deve imprimir uma linha de saída, contendo um valor real com quatro casas decimais, representando a aproximação obtida para $e$ com $M$ termos.

% \Notes
% Preste atenção aos tipos das variáveis, pois $F(90) \approx 2^{61}$.

\Example

\begin{example}
\exmp{
1
}{
1.0000
}%
\end{example}

\begin{example}
\exmp{
5
}{
2.7083
}%
\end{example}

\begin{example}
\exmp{
10
}{
2.7183
}%
\end{example}

% \textit{Problema baseado na questão ``A idade de dona Mônica'' da OBI2019 fase local.}

\end{problem}

\end{document}
